% ============================================================================
%  Theoretical Framework, part B: parametric models, soft constraints,
%  networks and operators. Input from main.tex right after ch2_theory.tex.
% ============================================================================

% ============================================================================
\section{SVI: three parameterisations, exact conversions, and the wings}\label{sec:svisec}
% ============================================================================

Everything so far has been a \emph{test}: given a surface, decide whether it is arbitrage-free.
Nothing so far tells us how to \emph{build} one. The parametric approach answers that by
restricting attention to a family small enough that the conditions of the previous sections can
be checked, and in the best case proved, once and for all in terms of the parameters.

SVI is that family. It has survived two decades of practitioner use for reasons that are worth
naming, because each of them is used somewhere in this thesis: its derivatives are available in
closed form, which is what allows us to validate automatic differentiation to machine precision
in Section~\ref{sec:deep}; it admits three equivalent parameterisations, one convenient for
calibration, one for the link to stochastic volatility, and one that traders can read off a
smile; and its wing behaviour matches the model-free bound of Lee exactly. This section develops
the three parameterisations and the exact maps between them. The maps matter more than they may
appear: they are the vehicle of the butterfly-repair algorithm of
Section~\ref{sec:parametric}, and round-tripping through them is what exposed an error in the
published inverse (Remark~\ref{rem:gjtypo}).

\subsection{The raw parameterisation and its derivatives}

The \emph{raw SVI} slice \cite{gatheral2004} at a fixed maturity $t$ is
\begin{equation}\label{eq:svi}
w(k;\chi_R)=a+b\Bigl(\rho\,(k-m)+\sqrt{(k-m)^2+\sigma^2}\Bigr),
\qquad \chi_R=(a,b,\rho,m,\sigma),
\end{equation}
with $b\ge0$, $|\rho|<1$, $\sigma>0$, and $a+b\sigma\sqrt{1-\rho^2}\ge0$ (which, as shown in
Lemma~\ref{lem:jw} below, is exactly nonnegativity of the minimum of $w$). Its derivatives are
closed-form:
\begin{equation}\label{eq:svideriv}
w'(k)=b\Bigl(\rho+\frac{k-m}{\sqrt{(k-m)^2+\sigma^2}}\Bigr),\qquad
w''(k)=\frac{b\,\sigma^2}{\bigl((k-m)^2+\sigma^2\bigr)^{3/2}}\;>\;0 .
\end{equation}
These closed forms are used in Section~\ref{sec:deep} to validate automatic differentiation to
machine precision. Note $w''>0$ everywhere: an SVI slice is strictly convex in $k$, yet convexity of $w$ in $k$ does not imply convexity of the
call in $K$, i.e.\ does not imply $g\ge0$.

\subsection{Wing asymptotics and Lee consistency}

From \eqref{eq:svi}, as $k\to\pm\infty$,
\begin{equation}\label{eq:sviwings}
w(k)=b\,(1\pm\rho)\,|k|+O(1),
\end{equation}
so the asymptotic total-variance slopes are $b(1+\rho)$ (right) and $b(1-\rho)$ (left). Lee's
moment formula \cite{lee2004} bounds the asymptotic slope of total variance in any
arbitrage-free slice: $\limsup_{k\to\pm\infty} w(k)/|k|\le 2$. Applying this to both wings gives
the \emph{wing condition}
\begin{equation}\label{eq:wingcond}
b\,(1+|\rho|)\;\le\;2 .
\end{equation}
with strict inequality required for condition (ii) of Lemma~\ref{lem:gjbutterfly}
($w(k)/k\to2$ is precisely $d_+\not\to-\infty$). Equation \eqref{eq:wingcond} constrains the
wings only; the interior condition $g\ge0$ has no known tractable characterisation in terms of
$\chi_R$ and must be checked (Section~\ref{sec:parametric} does so on a fine audit grid).

\subsection{The natural parameterisation and the corrected inverse map}

The \emph{natural} parameterisation writes the same family as a perturbation of a
time-rescaled Heston-consistent form:
\begin{equation}\label{eq:svinat}
w(k;\chi_N)=\Delta+\frac{\omega}{2}
\Bigl\{1+\zeta\rho\,(k-\mu)+\sqrt{\bigl(\zeta(k-\mu)+\rho\bigr)^2+1-\rho^2}\Bigr\},
\qquad \chi_N=(\Delta,\mu,\rho,\omega,\zeta).
\end{equation}

\begin{lemma}[Raw $\leftrightarrow$ natural, {cf.\ \cite[Lemma~3.2]{gatheraljacquier2014}}]
\label{lem:rawnat}
The parameterisations \eqref{eq:svi} and \eqref{eq:svinat} describe the same slice under
\begin{equation}\label{eq:nat2raw}
(a,b,\rho,m,\sigma)
=\Bigl(\Delta+\tfrac{\omega}{2}(1-\rho^2),\;\tfrac{\omega\zeta}{2},\;\rho,\;
\mu-\tfrac{\rho}{\zeta},\;\tfrac{\sqrt{1-\rho^2}}{\zeta}\Bigr),
\end{equation}
with inverse
\begin{equation}\label{eq:raw2nat}
(\Delta,\mu,\rho,\omega,\zeta)
=\Bigl(a- b\sigma\sqrt{1-\rho^2} ,\;
m+\tfrac{\rho\sigma}{\sqrt{1-\rho^2}},\;\rho,\;
\tfrac{2b\sigma}{\sqrt{1-\rho^2}},\;\tfrac{\sqrt{1-\rho^2}}{\sigma}\Bigr).
\end{equation}
\end{lemma}

\begin{proof}
Set $v:=\zeta(k-\mu)+\rho$. With $m=\mu-\rho/\zeta$ we have $k-m=v/\zeta$, so the raw form
reads
\[
a+b\Bigl(\rho\frac{v}{\zeta}+\sqrt{\frac{v^2}{\zeta^2}+\sigma^2}\Bigr)
=a+\frac{b}{\zeta}\Bigl(\rho v+\sqrt{v^2+\zeta^2\sigma^2}\Bigr).
\]
With $\sigma=\sqrt{1-\rho^2}/\zeta$, $\zeta^2\sigma^2=1-\rho^2$. The natural form expands as
\[
\Delta+\frac{\omega}{2}\bigl\{1+\rho(v-\rho)\bigr\}+\frac{\omega}{2}\sqrt{v^2+1-\rho^2}
=\Bigl(\Delta+\frac{\omega}{2}(1-\rho^2)\Bigr)
+\frac{\omega}{2}\Bigl(\rho v+\sqrt{v^2+1-\rho^2}\Bigr).
\]
Matching coefficients gives $b/\zeta=\omega/2$ and $a=\Delta+\omega(1-\rho^2)/2$, i.e.\
\eqref{eq:nat2raw}. Solving \eqref{eq:nat2raw} for $\chi_N$: $\zeta=\sqrt{1-\rho^2}/\sigma$,
$\omega=2b/\zeta=2b\sigma/\sqrt{1-\rho^2}$, $\mu=m+\rho/\zeta=m+\rho\sigma/\sqrt{1-\rho^2}$,
and
\[
\Delta=a-\frac{\omega}{2}(1-\rho^2)=a-\frac{2b\sigma}{2\sqrt{1-\rho^2}}(1-\rho^2)
=a-b\sigma\sqrt{1-\rho^2}.
\]
\end{proof}

\begin{remark}[A correction to the published inverse]\label{rem:gjtypo}
During the replication of Section~\ref{sec:parametric} we identified that the inverse map as
printed in \cite[Lemma~3.2]{gatheraljacquier2014} omits the $-b\sigma\sqrt{1-\rho^2}$
contribution to $\Delta$ (equivalently, a missing $-\alpha\sqrt{1-\rho^2}$ term in the paper's
notation). The derivation above fixes it; the error was caught by round-tripping $2{,}000$
random parameter draws through raw$\to$natural$\to$JW$\to$raw, which reaches $2.7\times10^{-10}$
maximum error only with the corrected term.
\end{remark}

\subsection{The Jump-Wings parameterisation}

The raw and natural parameterisations are convenient to compute with but opaque to read: a quintuple such as $(a,b,\rho,m,\sigma)$ conveys nothing directly about the shape of the smile it
encodes. The \emph{Jump-Wings} (JW) parameterisation \cite{gatheraljacquier2014} re-expresses
the \emph{same} slice through five quantities that a trader reads straight off the smile, each
corresponding to a visible geometric feature. At maturity $t$, with $w_t:=v_t\,t$ the ATM total
variance, these are the ATM variance $v_t$ (the level of the smile at the money), the ATM skew
$\psi_t$ (its slope at the money), the left and right wing slopes $p_t$ and $c_t$ (the asymptotic
steepness of the deep-OTM put and call wings, which carry the market's jump information), and the minimum implied variance $\tilde v_t$ (the level at the bottom of the smile,generally displaced from the money when $\psi_t\neq0$). 

\begin{lemma}[JW from raw, {\cite[eq.~(3.5)]{gatheraljacquier2014}}]\label{lem:jw}
For the raw slice \eqref{eq:svi} at maturity $t$,
\begin{equation}\label{eq:jw}
\begin{aligned}
v_t&=\frac{a+b\bigl(-\rho m+\sqrt{m^2+\sigma^2}\bigr)}{t},
&\psi_t&=\frac{1}{\sqrt{w_t}}\cdot\frac{b}{2}\Bigl(\rho-\frac{m}{\sqrt{m^2+\sigma^2}}\Bigr),\\
p_t&=\frac{b\,(1-\rho)}{\sqrt{w_t}},
&c_t&=\frac{b\,(1+\rho)}{\sqrt{w_t}},\qquad
\tilde v_t=\frac{a+b\sigma\sqrt{1-\rho^2}}{t}.
\end{aligned}
\end{equation}
\end{lemma}

\begin{proof}
$v_t t=w(0)$ follows by direct evaluation of \eqref{eq:svi} at $k=0$. $\psi_t$ is by
definition $w'(0)/(2\sqrt{w_t})$; insert \eqref{eq:svideriv} at $k=0$. $p_t$ and $c_t$ are the
normalised wing slopes $\lim_{k\to\mp\infty}|w'(k)|/\sqrt{w_t}$; by \eqref{eq:sviwings} these
are $b(1\mp\rho)/\sqrt{w_t}$ with the stated sign convention. For $\tilde v_t$: minimise
$u\mapsto\rho u+\sqrt{u^2+\sigma^2}$ ($u=k-m$). The first-order condition
$\rho+u/\sqrt{u^2+\sigma^2}=0$ gives $u^\star=-\rho\sigma/\sqrt{1-\rho^2}$, whence
$\sqrt{(u^\star)^2+\sigma^2}=\sigma/\sqrt{1-\rho^2}$ and the minimum value
$\rho u^\star+\sigma/\sqrt{1-\rho^2}=\sigma\sqrt{1-\rho^2}$; therefore
$\min_kw=a+b\sigma\sqrt{1-\rho^2}=\tilde v_t\,t$.
\end{proof}

\begin{lemma}[Raw from JW, {\cite[eq.~(3.6)]{gatheraljacquier2014}}]\label{lem:jwinv}
Given $(v_t,\psi_t,p_t,c_t,\tilde v_t)$ with $w_t=v_tt$, set
\begin{equation}\label{eq:jwinv}
b=\frac{\sqrt{w_t}}{2}\,(c_t+p_t),\qquad
\rho=1-\frac{p_t\sqrt{w_t}}{b}=\frac{c_t-p_t}{c_t+p_t},\qquad
\beta:=\rho-\frac{2\psi_t\sqrt{w_t}}{b},\qquad
\alpha:=\sgn(\beta)\sqrt{\frac1{\beta^2}-1},
\end{equation}
(requiring $|\beta|\le1$ for a real solution), and then
\begin{equation}\label{eq:jwinv2}
m=\frac{(v_t-\tilde v_t)\,t}
{b\,\bigl\{-\rho+\sgn(\alpha)\sqrt{1+\alpha^2}-\alpha\sqrt{1-\rho^2}\bigr\}},\qquad
\sigma=\alpha\,m.
\end{equation}
The degenerate case $m=0$ occurs exactly when $v_t=\tilde v_t$, i.e.\ when the minimum of the
slice sits at the money, so that the numerator of \eqref{eq:jwinv2} vanishes and the formula
$\sigma=\alpha m$ becomes indeterminate; it is handled separately by
$\sigma=\bigl((v_t-\tilde v_t)\,t\bigr)/\bigl(b\,(1-\sqrt{1-\rho^2})\bigr)$. Finally,
$a=\tilde v_t\,t-b\sigma\sqrt{1-\rho^2}$.
\end{lemma}

\begin{proof}[Proof sketch]
Adding and subtracting the wing equations of \eqref{eq:jw} gives $b$ and $\rho$ directly. The
skew equation then determines $m/\sqrt{m^2+\sigma^2}=\rho-2\psi_t\sqrt{w_t}/b=\beta$, i.e.\
$\sigma=\alpha m$ with $\alpha$ as stated (the sign bookkeeping selects the branch). Finally
$v_t-\tilde v_t$ measures the gap between ATM and minimum variance, which in raw parameters is
$b\{-\rho m+\sqrt{m^2+\sigma^2}-\sigma\sqrt{1-\rho^2}\}$; substituting $\sigma=\alpha m$ and
solving for $m$ yields \eqref{eq:jwinv2}, and $a$ follows from
$\tilde v_tt=a+b\sigma\sqrt{1-\rho^2}$ (Lemma~\ref{lem:jw}).
\end{proof}

\subsection{The Vogt smile: arbitrage invisible to the eye}

\begin{example}[{Vogt smile, \cite[Ex.~3.1]{gatheraljacquier2014}}]\label{ex:vogt}
Consider the raw SVI slice at $t=1$ with parameters
$$
(a,\,b,\,m,\,\rho,\,\sigma)=(-0.0410,\;0.1331,\;0.3586,\;0.3060,\;0.4153).
$$

This slice looks entirely healthy: its total variance is positive everywhere, it is strictly
convex in $k$, and it satisfies the wing condition \eqref{eq:wingcond} with a wide margin,
$b(1+|\rho|)=0.174\ll2$. To the naked eye it is an ordinary smile. Yet its Durrleman function
dips below zero, $\min_k g=-0.0329$ (attained near $k\approx0.88$), so by
Proposition~\ref{prop:density} the implied density is negative on an interval: a genuine
butterfly arbitrage that no visual inspection would reveal. Figure~\ref{fig:vogt} reproduces
the exhibit from the closed forms \eqref{eq:svideriv}, and Section~\ref{sec:parametric}
replicates the paper's Jump--Wings quintuple for this slice digit for digit, together with its
repair recipe. The Vogt smile is the thesis's first illustration of a principle that recurs
throughout: visual smoothness certifies nothing, and arbitrage must be checked, not eyeballed.
\end{example}

\thfig{figures/fig_vogt.pdf}{1.0}{The Vogt SVI slice of Example~\ref{ex:vogt}, computed from the
closed forms \eqref{eq:svideriv}--\eqref{eq:g}. Left: the total variance $w(k)$ is positive,
smooth and convex in $k$, an entirely ordinary-looking smile. Centre: the Durrleman function
$g(k)$ nonetheless dips below zero, $\min_k g=-0.0329$ near $k\approx0.88$. Right: by
\eqref{eq:density} the implied density $q(k)$ is therefore negative on an interval (shaded), a
genuine butterfly arbitrage that no visual inspection of the left panel would reveal.}{fig:vogt}

\thfig{figures/fig_svi_wings.pdf}{1.0}{The same SVI slice read two ways, from the closed forms
\eqref{eq:svi}--\eqref{eq:sviwings}. Left: far from the money the slice is asymptotically linear,
with left-wing slope $b(1-\rho)$ and right-wing slope $b(1+\rho)$ (dashed); Lee's bound caps the
steeper of the two at $2$, which is the wing condition \eqref{eq:wingcond}. Right: two of the
five Jump--Wings quantities of Lemma~\ref{lem:jw} located directly on the curve, the ATM variance
$v_t\,t$ (dot, at $k=0$) and the minimum variance $\tilde v_t\,t$ (dotted line); the dot sits
above the trough because the skew displaces the minimum away from the money.}{fig:sviwings}

% ============================================================================
\section{SSVI and the Gatheral--Jacquier certificates}\label{sec:ssvisec}
% ============================================================================

\subsection{The surface}

SVI describes one maturity at a time. Fitting each maturity independently, as
Section~\ref{sec:parametric} does, yields excellent per-slice accuracy and no guarantee whatever
that adjacent slices are mutually consistent: nothing in a per-slice fit prevents two slices
from crossing, which by Section~\ref{sec:calendarsec} is calendar arbitrage. SSVI resolves this
by tying every slice to a single term structure, so that the calendar direction is controlled by
construction rather than repaired afterwards. The price is rigidity, and
Section~\ref{sec:parametric} measures it precisely.

What makes SSVI valuable for this thesis is not only the construction but what Gatheral and
Jacquier prove about it: closed-form conditions on the parameters under which the surface is
provably free of both arbitrages. That is a \emph{certificate}, not an audit, and it is the only
such object in the thesis. It is what the deep smoother of Section~\ref{sec:deep} uses as a
prior, and it is why the ablation finding that the prior does the protective work rests on a
theorem rather than on luck.

The \emph{SSVI} (surface SVI) parameterisation \cite{gatheraljacquier2014} promotes SVI to a
full surface driven by its ATM total-variance term structure
$\theta_\tau:=w(0,\tau)$:
\begin{equation}\label{eq:ssvi}
w(k,\theta_\tau)=\frac{\theta_\tau}{2}\Bigl(1+\rho\,\varphi(\theta_\tau)\,k
+\sqrt{\bigl(\varphi(\theta_\tau)k+\rho\bigr)^2+1-\rho^2}\Bigr),
\end{equation}
where $\varphi:(0,\infty)\to(0,\infty)$ is a smooth curvature function and $|\rho|<1$.
Evaluating at $k=0$ confirms $w(0,\theta)=\tfrac\theta2(1+\sqrt{\rho^2+1-\rho^2})=\theta$: the
term structure is matched by construction. In this thesis we use the \emph{power law}
$\varphi(\theta)=\eta\theta^{-\gamma}$ with a smooth increasing term structure
$\theta_\tau=\alpha\tau^{\beta}$ ($\eta,\alpha,\beta>0$, $\gamma\in(0,1)$).

Each SSVI slice is itself a raw SVI slice; this observation, proved next, transports the SVI
machinery (derivatives, $g$-audits, JW readings) to SSVI and is the first step of the
butterfly certificate.

\begin{lemma}[SSVI slices are SVI slices]\label{lem:ssvi2svi}
For fixed $\theta>0$, the slice $k\mapsto w(k,\theta)$ of \eqref{eq:ssvi} equals the raw SVI
slice \eqref{eq:svi} with parameters
\begin{equation}\label{eq:ssviraw}
(a,b,\rho,m,\sigma)
=\Bigl(\frac{\theta}{2}(1-\rho^2),\;\frac{\theta\varphi(\theta)}{2},\;\rho,\;
-\frac{\rho}{\varphi(\theta)},\;\frac{\sqrt{1-\rho^2}}{\varphi(\theta)}\Bigr).
\end{equation}
\end{lemma}

\begin{proof}
With the candidate parameters, $k-m=k+\rho/\varphi$, so
\[
b\,\rho(k-m)=\frac{\theta\varphi}{2}\,\rho\Bigl(k+\frac{\rho}{\varphi}\Bigr)
=\frac{\theta}{2}\bigl(\rho\varphi k+\rho^2\bigr),\qquad
b\sqrt{(k-m)^2+\sigma^2}
=\frac{\theta}{2}\sqrt{(\varphi k+\rho)^2+1-\rho^2},
\]
and adding $a=\theta(1-\rho^2)/2$ reproduces \eqref{eq:ssvi} exactly.
\end{proof}

\subsection{The calendar certificate}

\begin{proposition}[{Calendar, \cite[Thm.~4.1]{gatheraljacquier2014}}]\label{thm:gjcal}
The SSVI surface \eqref{eq:ssvi} is free of calendar arbitrage if and only if
\begin{enumerate}[label=(\roman*)]
\item $\tau\mapsto\theta_\tau$ is non-decreasing, and
\item $0\le\partial_\theta\bigl(\theta\varphi(\theta)\bigr)\le
\dfrac{1}{\rho^{2}}\Bigl(1+\sqrt{1-\rho^{2}}\Bigr)\varphi(\theta)$ for all $\theta>0$.
\end{enumerate}
\end{proposition}

\begin{proof}
By Proposition~\ref{prop:calendar} and the chain rule
$\partial_\tau w=\partial_\theta w\cdot\partial_\tau\theta_\tau$, calendar-freeness is
equivalent to (i) together with $\partial_\theta w(k,\theta)\ge0$ for all $(k,\theta)$. Write
$\psi(\theta):=\theta\varphi(\theta)$ and multiply \eqref{eq:ssvi} by $2$:
\[
2w(k,\theta)=\theta+\rho\,\psi(\theta)\,k
+\sqrt{\bigl(\psi(\theta)k+\rho\theta\bigr)^2+(1-\rho^2)\theta^2}
\]
 
Differentiating in $\theta$, with
$R:=\sqrt{(\psi k+\rho\theta)^2+(1-\rho^2)\theta^2}$, we obtain:
\begin{equation}\label{eq:dthetaw}
2\,\partial_\theta w
=1+\rho\,\psi'(\theta)\,k
+\frac{(\psi k+\rho\theta)\bigl(\psi'(\theta)k+\rho\bigr)+(1-\rho^2)\theta}{R}.
\end{equation}

\emph{Necessity of $\psi'\ge0$.} As $k\to+\infty$, $R\sim\psi k$ and
\eqref{eq:dthetaw} $\sim\rho\psi'k+\psi'k=(1+\rho)\,\psi'\,k$; as $k\to-\infty$,
$R\sim\psi|k|$ and \eqref{eq:dthetaw} $\sim\rho\psi'k+\psi'|k|=(1-\rho)\,\psi'\,|k|$. Since
$1\pm\rho>0$, nonnegativity of $\partial_\theta w$ in both wings forces $\psi'(\theta)\ge0$;
this is the left inequality of (ii).


\emph{The sharp upper bound.} Assume $\psi'\ge0$ and set
$x:=\psi(\theta)k+\rho\theta$, which defines a bijection in $k$ for fixed $\theta$.
A direct computation reduces the problem to minimising a one-variable function of $x$.
The minimiser is obtained by setting the $x$-derivative of \eqref{eq:dthetaw} to zero,
which yields a quadratic equation. Requiring the resulting minimum to be nonnegative shows
that $\inf_k\partial_\theta w\ge0$ holds precisely when


\[
\theta\,\frac{\psi'(\theta)}{\psi(\theta)}\;\le\;\frac{1+\sqrt{1-\rho^2}}{\rho^2},
\quad\text{i.e.}\quad
\psi'(\theta)\le\frac{1}{\rho^2}\bigl(1+\sqrt{1-\rho^2}\bigr)\varphi(\theta),
\]
the right inequality of (ii); the bound is attained (the constant
$(1+\sqrt{1-\rho^2})/\rho^2$ is sharp), so (ii) is also necessary. The complete algebra of
the minimisation is in the proof of \cite[Thm.~4.1]{gatheraljacquier2014}; the boundary cases
$\rho=\pm1$ are treated there by a separate limiting argument, and $\rho=0$ makes the upper
bound vacuous ($+\infty$), consistent with \eqref{eq:dthetaw} being manifestly nonnegative
when $\rho=0$ and $\psi'\ge0$.
\end{proof}

\begin{example}[Power law satisfies (ii) automatically]\label{ex:powerlawcal}
For $\varphi(\theta)=\eta\theta^{-\gamma}$, $\psi(\theta)=\eta\theta^{1-\gamma}$ and
$\psi'(\theta)=(1-\gamma)\eta\theta^{-\gamma}=(1-\gamma)\varphi(\theta)$. Condition (ii) reads
$0\le1-\gamma\le\kappa(\rho):=(1+\sqrt{1-\rho^2})/\rho^2$. Since
$\kappa(\rho)\ge1$ for all $|\rho|\le1$ (with equality only at $\rho=\pm1$) and
$1-\gamma\in(0,1)$, the condition holds for every $\gamma\in(0,1)$ and every $\rho$: with a
non-decreasing $\theta_\tau$ (e.g.\ $\alpha\tau^\beta$), the power-law SSVI is calendar-free
for free.
\end{example}

\subsection{The butterfly certificate}

\begin{proposition}[{Butterfly, \cite[Thm.~4.2]{gatheraljacquier2014}}]\label{thm:gjbfly}
The SSVI surface is free of butterfly arbitrage if, for all $\theta>0$,
\begin{equation}\label{eq:gjconds}
\theta\varphi(\theta)\,(1+|\rho|)<4
\qquad\text{and}\qquad
\theta\varphi(\theta)^{2}\,(1+|\rho|)\le4 .
\end{equation}
Together with Proposition~\ref{thm:gjcal} this yields freedom from \emph{static} arbitrage
(\cite[Cor.~4.1]{gatheraljacquier2014}).
\end{proposition}

\begin{proof}
Fix $\theta$ and use Lemma~\ref{lem:ssvi2svi}: the slice is raw SVI with
$b=\theta\varphi/2$ and correlation $\rho$.

\begin{enumerate}[label=\textup{(\roman*)}, leftmargin=*]

\item \textbf{Step 1: The first condition is the wing condition.}

By \eqref{eq:sviwings}, the wing slopes are $b(1\pm\rho)$, so
\[
b(1+|\rho|)
=\frac{\theta\varphi(\theta)}{2}(1+|\rho|)<2
\iff
\theta\varphi(\theta)(1+|\rho|)<4.
\]
Thus, the first condition in \eqref{eq:gjconds} is exactly the strict Lee wing bound
\eqref{eq:wingcond}. It yields condition~(ii) of Lemma~\ref{lem:gjbutterfly}, namely
$d_+\to-\infty$ and the vanishing of call prices at extreme strikes.

\item \textbf{Step 2: The second condition ensures $g\ge0$.}

Substituting the slice parameters \eqref{eq:ssviraw} into the Durrleman function
\eqref{eq:g} and setting $z:=\varphi(\theta)k$ reduces $g\ge0$ on $\R$ to the
nonnegativity of an explicit rational-algebraic function of $z$. Its coefficients depend
only on $\theta\varphi^2$ and $\rho$. As shown in
\cite[Thm.~4.2]{gatheraljacquier2014}, the conditions
\[
\theta\varphi(\theta)^2(1+|\rho|)\le4
\]
and the wing bound from Step~1 ensure that this function is nonnegative for every $z$.

The quantity $\theta\varphi^2$ represents the natural curvature scale of the slice. Indeed,
from \eqref{eq:ssviraw},
\[
\sigma=\frac{\sqrt{1-\rho^2}}{\varphi},
\qquad
w''(m)=\frac{b}{\sigma}
=\frac{\theta\varphi^2}{2\sqrt{1-\rho^2}}.
\]
The second condition therefore controls the curvature near the money relative to the level
$\theta$, which is precisely the combination through which the terms amplified by $1/w$ in
\eqref{eq:g} may make $g$ negative.\medskip

The full algebraic computation of Step 2 is carried out line by line in the proof of
\cite[Thm.~4.2, pp.~14--15]{gatheraljacquier2014}; we do not reproduce it here. In this thesis
it is instead encoded directly: Propositions~\ref{thm:gjcal}--\ref{thm:gjbfly} are implemented as
\emph{executable certificates} and verified against the paper's own condition values
(Section~\ref{sec:parametric}).
\end{enumerate}
\end{proof}

\begin{example}[$\varphi$ families]\label{ex:phifamilies}
\emph{Heston-like:} $\varphi(\theta)=\frac{1}{\lambda\theta}
\bigl(1-\frac{1-e^{-\lambda\theta}}{\lambda\theta}\bigr)$, $\lambda>0$, satisfies
\eqref{eq:gjconds} for all $\theta$ iff $\lambda\ge(1+|\rho|)/4$, and condition (ii) of
Proposition~\ref{thm:gjcal} always. \emph{Power law:} calendar-free by
Example~\ref{ex:powerlawcal}, but see Remark~\ref{rem:gj44}. \emph{Modified power law:}
$\varphi(\theta)=\eta\,\theta^{-\gamma}(1+\theta)^{\gamma-1}$ with $\gamma\in(0,\tfrac12]$ and
$\eta(1+|\rho|)\le2$ satisfies \eqref{eq:gjconds} for \emph{all} $\theta>0$: a globally
certified family. All three are implemented and verified in NB02 with the paper's published
condition values.
\end{example}

\begin{remark}[{Why the plain power law cannot be certified globally;
\cite[Rem.~4.4]{gatheraljacquier2014}}]\label{rem:gj44}
For $\varphi(\theta)=\eta\theta^{-\gamma}$: $\theta\varphi=\eta\theta^{1-\gamma}\to\infty$ as
$\theta\to\infty$ (violating the first condition of \eqref{eq:gjconds} at large $\theta$ for
any $\gamma\in(0,1)$), and if $\gamma>\tfrac12$,
$\theta\varphi^2=\eta^2\theta^{1-2\gamma}\to\infty$ as $\theta\downarrow0$ (violating the
second at small $\theta$). Hence no choice of $(\eta,\gamma)$ certifies the plain power law on
all of $(0,\infty)$. The next lemma shows this is the wrong question for a smoother: what a
smoother evaluates is a \emph{compact working range} of $\theta$, and on a compact range the
certificate reduces to two endpoint checks.
\end{remark}

\subsection{Endpoint certification on the working range}

\begin{theorem}[Endpoint certification]\label{lem:range}
Let $\varphi(\theta)=\eta\theta^{-\gamma}$ with $\eta>0$, $\gamma\in(0,1)$, and let
$0<\theta_{\min}\le\theta\le\theta_{\max}$ be the range attained by
$\theta_\tau=\alpha\tau^{\beta}$ on the working maturities $[\tau_{\min},\tau_{\max}]$
($\alpha,\beta>0$). Then
\[
\sup_{\theta\in[\theta_{\min},\theta_{\max}]}\theta\varphi(\theta)
=\eta\,\theta_{\max}^{1-\gamma},
\qquad
\sup_{\theta\in[\theta_{\min},\theta_{\max}]}\theta\varphi(\theta)^{2}
=\eta^{2}\max\bigl(\theta_{\max}^{\,1-2\gamma},\,\theta_{\min}^{\,1-2\gamma}\bigr).
\]
Consequently, checking \eqref{eq:gjconds} at the two endpoints certifies the SSVI surface free
of static arbitrage on the entire working domain $[\tau_{\min},\tau_{\max}]\times\R$.
\end{theorem}

\begin{proof}
$\theta\varphi(\theta)=\eta\theta^{1-\gamma}$ is increasing for $\gamma<1$, so its supremum is
at $\theta_{\max}$. $\theta\varphi(\theta)^2=\eta^2\theta^{1-2\gamma}$ is increasing iff
$\gamma\le\tfrac12$ (supremum at $\theta_{\max}$) and decreasing iff $\gamma>\tfrac12$
(supremum at $\theta_{\min}$); the displayed maximum covers both. Monotonicity of
$\theta_\tau=\alpha\tau^\beta$ maps $[\tau_{\min},\tau_{\max}]$ onto
$[\theta_{\min},\theta_{\max}]$, so the working $\theta$-range is exactly this interval;
conditions \eqref{eq:gjconds} on the range then follow from the two endpoint evaluations, and
Proposition~\ref{thm:gjbfly} plus Example~\ref{ex:powerlawcal} plus Proposition~\ref{thm:gjcal}(i)
give static-arbitrage-freeness on $[\tau_{\min},\tau_{\max}]\times\R$.
\end{proof}

\thfig{figures/fig_phi_certificate.pdf}{1.0}{Endpoint certification
(Theorem~\ref{lem:range}): $\theta\varphi(\theta)$ is monotone increasing (left), so its sup on
a working range sits at $\theta_{\max}$; $\theta\varphi(\theta)^2$ (right, log scale) is
monotone with direction set by $\gamma\lessgtr\tfrac12$, so its sup sits at one of the two
endpoints. Checking the bound $4/(1+|\rho|)$ (dashed red) at the endpoints certifies the whole
range (dotted).}{fig:phicert}

% \begin{remark}\label{rem:cert}
% Theorem~\ref{lem:range} converts GJ from an \emph{audit} into a \emph{constraint}: our prior
% calibrator barriers \eqref{eq:gjconds} at the endpoints during the search and hard-projects
% $\eta$ afterwards, then prints a certificate. Anything below $\tau_{\min}$ remains uncertified
% extrapolation and is flagged as such downstream (the Monte-Carlo time grids of
% Chapter~\ref{ch:downstream} are floored at $\tau_{\min}$). An earlier version of the
% calibrator clipped $\gamma\le\tfrac12$ to certify on $(0,\theta_{\max}]$; on SPX data the fit
% prefers $\gamma>\tfrac12$ and the clip cost $\approx0.5\,\vp$ of prior accuracy---the endpoint
% form recovers it at no loss of guarantee.
% \end{remark}

\newpage

\subsection{Interpolation, extrapolation, and crossedness}
Two further Gatheral--Jacquier results are used by the parametric benchmark. First, between two
calibrated slices $w_1,w_2$ at ATM total variances $\theta_1<\theta_2$, butterfly-free and
non-crossing, \cite[Lem.~5.1]{gatheraljacquier2014} construct an \emph{arbitrage-free
interpolation} at any intermediate $\theta_t\in[\theta_1,\theta_2]$. The interpolation is
performed in \emph{price} space, not directly on $w$: with the $\sqrt{\theta}$-weight
\begin{equation}\label{eq:interp}
\alpha=\frac{\sqrt{\theta_2}-\sqrt{\theta_t}}{\sqrt{\theta_2}-\sqrt{\theta_1}}\in[0,1],
\end{equation}
one forms the convex combination of the two normalised call slices and inverts it back to total
variance through the Black map \eqref{eq:blacknorm},
\begin{equation}\label{eq:interp2}
w(k,\theta_t)=c^{-1}\!\Bigl(k,\;\alpha\,c\bigl(k,w_1(k)\bigr)+(1-\alpha)\,c\bigl(k,w_2(k)\bigr)\Bigr),
\end{equation}
where $c^{-1}(k,\cdot)$ denotes Black inversion at fixed $k$ (unique since $c_w>0$,
Lemma~\ref{lem:blackpartials}). Convexity of the price in each slice is preserved by the convex
combination, so the interpolated slice is butterfly-free by construction; beyond the last
calibrated maturity we extrapolate by the dominating shift $w_{\mathrm{ext}}(k)=w_2(k)+(\theta_{\mathrm{far}}-\theta_2)$,
which raises the whole slice and therefore introduces no calendar crossing. We implement both and
verify them numerically: the interpolated slice attains $\min_k g=0.52$, and the extrapolated
slice dominates the last calibrated slice everywhere, so no static arbitrage is introduced by
either operation. Second, when fitted slices \emph{do} cross, the crossing is quantified by the
\emph{crossedness} of \cite[Def.~5.1]{gatheraljacquier2014}: for adjacent slices $\chi_1,\chi_2$
with intersection points $\tilde k_1,\dots,\tilde k_n$, the crossedness is the maximum over these
points of $\max\bigl(0,\,w(\tilde k;\chi_1)-w(\tilde k;\chi_2)\bigr)$; a penalty on total
crossedness is the calendar-repair device of Section~\ref{sec:parametric}.

\thfig{figures/fig_ssvi.pdf}{1.0}{The SSVI surface \eqref{eq:ssvi} with power-law
$\varphi$ and $\theta_\tau=\alpha\tau^\beta$ ($\rho=-0.7$, $\eta=0.8$, $\gamma=0.45$,
$\alpha=0.10$, $\beta=0.9$; parameters satisfy \eqref{eq:gjconds} on the plotted range).
Left: total-variance slices never cross (Proposition~\ref{thm:gjcal}). Right: the corresponding
implied-volatility surface.}{fig:ssvi}
